% TOP_PROBLEM: {"id": 3032, "title": "Waring rank of determinant", "category_id": 4, "category": {"id": 4, "name": "algebra", "display_name": "Algebra", "description": "Group theory, ring theory, field theory, and algebraic structures.", "slug": "algebra", "order_index": 4, "created_at": "2026-07-31T15:26:25.671Z"}, "status": "open", "problem_number": "OPG-60031", "set_id": 12}
% FIRST_POSTED: 2026-10-01
% ATTEMPT_STATUS: unresolved
% UnsolvedMath ID 3032; source catalogue code OPG-60031
% !TeX program = lualatex
% Research attempt by GPT-6 Astra Ultra; no claim of a new complete solution.
\documentclass[11pt,a4paper]{article}
\usepackage[margin=25mm]{geometry}
\usepackage{fontspec}
\setmainfont{lmroman10-regular.otf}[BoldFont=lmroman10-bold.otf,
 ItalicFont=lmroman10-italic.otf,BoldItalicFont=lmroman10-bolditalic.otf]
\usepackage{amsmath,amssymb,mathtools}
\usepackage{unicode-math}
\setmathfont{latinmodern-math.otf}
\AtBeginDocument{\let\setminus\smallsetminus}
\usepackage{xurl,hyperref}
\hypersetup{colorlinks=true,urlcolor=blue}
\setcounter{secnumdepth}{0}
\setlength{\parindent}{0pt}
\setlength{\parskip}{5pt}
\setlength{\emergencystretch}{3em}
\begin{document}
\section*{Waring rank of determinant}\label{3032}
{\small UnsolvedMath ID 3032 · OPG-60031 · Algebra · OpenGarden}
\subsection{Definitions and mathematical statement}
Fix an integer $d\ge1$ and work in the polynomial ring
$R=\mathbb C[x_{ij}:1\le i,j\le d]$, whose $d^2$ variables are independent.
The generic determinant is the homogeneous degree-$d$ polynomial
\[
 D_d=\det(x_{ij})=
 \sum_{\sigma\in S_d}\operatorname{sgn}(\sigma)
       \prod_{i=1}^d x_{i,\sigma(i)},
\]
where $S_d$ is the group of permutations of $\{1,\ldots,d\}$.
A homogeneous linear form is an expression
$\ell=\sum_{i,j}a_{ij}x_{ij}$ with coefficients $a_{ij}\in\mathbb C$ and no
constant term.

The complex Waring rank of a nonzero homogeneous degree-$d$ form $F$ is
\[
 R_{\mathrm W}(F)=\min\left\{r\ge1:
 F=\sum_{\nu=1}^r\ell_\nu^d\text{ identically in }R\right\}.
\]
Allowing scalar coefficients in front of the powers gives the same number,
because each complex scalar has a $d$th root and can be absorbed into its
linear form. The forms may involve all entries of the matrix and need not
satisfy any symmetry or use disjoint variables.

Determine $R_{\mathrm W}(D_d)$ as a function of $d$. An exact answer consists
of a power-sum expression achieving the proposed value and a lower bound
excluding every expression with fewer summands. The equality is polynomial
identity on all matrices, including singular matrices. It is not enough to
express the determinant as products of unrelated linear forms, or as a
limit of shorter power sums: these concern other rank notions. In particular,
border rank and unsymmetrized tensor rank are not the invariant asked for.
\subsection{Short English statement}
What is the smallest number of complex $d$th powers of linear forms whose sum is the generic $d$ by $d$ determinant?
\subsection{Research attempt}
\paragraph{Scope and process.}
This is a GPT-6 Astra Ultra research attempt dated 1 October 2026, using
primary-source browsing and elementary mathematical checks. The canonical
definition above is retained. Results below are restricted results or
reductions, not a claim of a new solution to the entire catalogue question.
No formal proof assistant or external mathematical referee checked this text.


\paragraph{Literature and conventions.}
The rank is complex Waring rank, with literal sums of powers as in the
definition. The primary paper [R1] develops stronger lower bounds for
determinants using invariant forms. We use only derivative spaces and
polarization, proved below; the resulting estimates are benchmarks and
are not advertised as the best known bounds.

\paragraph{Derivative-space lower bound.}
For a degree-$d$ polynomial $F$, let $V_k(F)$ be the span of all its
partial derivatives of total order $k$, where $0\le k\le d$.
If $F=\sum_{\nu=1}^r\ell_\nu^d$, every such derivative is a scalar
multiple of some $\ell_\nu^{d-k}$, summed over $\nu$. Hence
$\dim V_k(F)\le r$. This uses the actual decomposition and gives
a lower bound for the rank requested in the problem.

For $D_d$, a nonzero order-$k$ derivative differentiates distinct
rows and distinct columns, and yields, up to sign, a complementary
$(d-k)$-minor. Every such minor occurs: pair the omitted rows and
columns in any order and differentiate in those $k$ entries.
Derivatives repeating a variable, row or column vanish by the
multilinearity of each determinant monomial. Distinct minors of the
same size are linearly independent, because each monomial determines
its row set and column set and hence occurs in just one of these minors.
This remains valid for the unique empty minor when $k=d$.
Consequently
\[
 \dim V_k(D_d)=\binom dk^2,\qquad
 R_{\mathrm W}(D_d)\ge\binom d{\lfloor d/2\rfloor}^{\!2}.
\]

\paragraph{An explicit general upper bound.}
For independent variables $y_1,\ldots,y_d$, character cancellation on
the sign cube gives
\[
 y_1\cdots y_d=
 \frac{1}{2^d d!}\sum_{\epsilon\in\{-1,1\}^d}
 \left(\prod_i\epsilon_i\right)
 \left(\sum_i\epsilon_i y_i\right)^d.
\]
To check this, expand the powers. The sign sum of a monomial with
exponents $a_i$ is zero unless every $a_i$ is odd. Since their sum
is $d$, this forces every $a_i=1$. Its multinomial coefficient is
$d!$, so the surviving coefficient is one. The summands indexed by
$\epsilon$ and $-\epsilon$ are identical, including the prefactor,
so only $2^{d-1}$ powers are needed. Apply this separately to the
$d!$ signed monomials in the determinant and absorb each nonzero
complex coefficient into a $d$th root. This proves
\[
 R_{\mathrm W}(D_d)\le d!\,2^{d-1}.
\]
Possible cancellations between these powers can only improve this bound.

\paragraph{Exact low-dimensional checks.}
For $d=1$ the rank is one. For $d=2$, write the entries as $a,b,c,e$;
then
\[
 ae-bc=\frac14\bigl((a+e)^2-(a-e)^2
                     -(b+c)^2+(b-c)^2\bigr).
\]
This is a four-square complex decomposition after absorbing signs and
the factor $1/4$. The first derivatives are $e,-c,-b,a$, which are
linearly independent, so fewer than four powers are impossible. Thus
$R_{\mathrm W}(D_2)=4$. Over $\mathbb R$, requiring sums of squares
with no signed coefficients would be a different problem, so the
complex-field hypothesis is essential to this verification.

\paragraph{Where the attack stops.}
Already for $d=3$ these elementary arguments give only $9\le R\le24$.
The derivative dimension records the span generated by potential powers,
but does not show that a spanning collection can actually realize the
determinant with that many powers. Conversely, treating permutation
monomials independently wastes cancellations. Neither equality in the
lower bound nor optimality of the upper construction has been proved.
Higher-dimensional apolar or representation-theoretic restrictions,
such as those studied in [R1], are additional information, not a
consequence of the displayed dimensions.

\paragraph{Final status.}
\textbf{UNRESOLVED} as a function of arbitrary $d$. The exact values
for $d=1,2$ and general elementary two-sided bounds are established;
no new exact higher-dimensional rank is claimed.

\subsection{Sources}
\begin{itemize}
\item[S1] ULAM.AI and the UnsolvedMath Contributors, supplied problems.json, record 3032 (OPG-60031), OpenGarden collection. Catalogue identity and source transcription; CC BY 4.0.
\url{https://huggingface.co/datasets/ulamai/UnsolvedMath}.
\item[S2] Open Problem Garden, Waring rank of determinant. Original problem and discussion; the mathematical formulation below is expanded from the supplied source record.
\url{http://www.openproblemgarden.org/op/waring_rank_of_determinant}.

\item[R1] Harm Derksen and Zach Teitler, \emph{Lower bound for ranks of
invariant forms}, primary research article, author-hosted copy.
\url{https://sites.lsa.umich.edu/hderksen/wp-content/uploads/sites/614/2018/05/A.I.a.43.pdf}.

\end{itemize}
\end{document}
